\documentclass[11pt, a4paper]{article}

% --- Packages ---
\usepackage[utf8]{inputenc}
\usepackage{amsmath}
\usepackage{graphicx}
\usepackage{geometry}
\usepackage{float}
\usepackage{caption}
\usepackage{booktabs}

\geometry{margin=1in}

% --- Title Information ---
\title{{\bf MLDS HOMEWORK 2} \\
(MODEL EVALUATION)}
\author{Denis Nedić}
\date{\today}

\begin{document}

\maketitle
\vspace{-1cm} % Pulls the first section up slightly to save space

\section{\bf{Part 1} Model evaluation and comparison}
\subsection{Methodology and results}
To predict basketball shot types, three models were evaluated: \textbf{Baseline classifier} (predicting relative class frequencies), \textbf{Logistic Regression}, and a \textbf{Decision Tree}.
Because Decision Trees are highly sensitive to overfitting, the maximum depth parameter was tuned using two distinct methodologies: optimizing directly on the training fold
and employing nested cross-validation. Performance was estimated using standard cross-validation, with Classification Accuracy and Log-Score serving as the primary evaluation metrics.
\begin{table}[htpb]
    \centering
    \caption{Cross-Validation Performance Metrics (Mean $\pm$ Standard Deviation)}
    \label{tab:model_comparison}
    \begin{tabular}{lcc}
        \toprule
        \textbf{Model} & \textbf{Accuracy} & \textbf{Log-Score} \\
        \midrule
        Baseline Classifier & $0.6081 \pm 0.0132$ & $1.1658 \pm 0.0282$ \\
        Logistic Regression & $0.7361 \pm 0.0245$ & $\mathbf{0.6729 \pm 0.0545}$ \\
        Decision Tree (Train-Opt) & $0.7034 \pm 0.0243$ & $10.0925 \pm 0.8872$ \\
        Decision Tree (Nested CV) & $\mathbf{0.7486 \pm 0.0230}$ & $1.0858 \pm 0.2290$ \\
        \bottomrule
    \end{tabular}
\end{table}
Table \ref{tab:model_comparison} summarizes the cross-validation results, reporting the mean and standard deviation across folds. The Decision Tree evaluated via Nested CV, which utilized 10-fold
cross-validation for both the inner parameter tuning loop and the outer evaluation loop. It achieved the highest overall classification accuracy, while Logistic Regression achieved the most optimal 
(lowest) log-score. Most notably, optimizing the Decision Tree directly on the training fold resulted in slightly lower accuracy but a catastrophic, highly volatile log-score. This severe penalty 
occurs because standard training fold optimization leads to massive overfitting; the tree creates overly deep, pure leaves that output overconfident, incorrect predictions (probabilities near 1.0 or 0.0), 
which are heavily penalized by log-loss. Employing nested cross-validation provided a much more stable and unbiased evaluation by strictly separating the depth-tuning process from the validation data, mitigating this overconfidence and significantly improving both metrics.
\section{{\bf Part 2}}
\subsection{Error dependency}
\begin{figure}[htpb]
    \centering
    \includegraphics[width=0.7\linewidth]{HW2/error.png} 
    \caption{Actual Error Rate and Model Uncertainty by 5-Degree Angle Bins.}
    \label{fig:angle_plots}
\end{figure}
To address the suspicion that prediction error depends on the shot angle, the continuous angle variable was discretized into 5-degree bins. Visual inspection of the binned error rates (Figure \ref{fig:angle_plots}) revealed a complex, wave-like distribution for both the Logistic Regression and Decision Tree models.

To formally evaluate this relationship, a Pearson's Chi-squared test of independence was utilized. This methodology evaluates whether prediction mistakes are disproportionately concentrated in specific 
angle sectors by treating the bins categorically, requiring no underlying assumptions about the mathematical shape of the data. The test yielded highly statistically significant p-values for both models (p$<$0.001). 
Therefore, we firmly confirm the initial suspicion: the model's error rates are heavily dependent on the shot angle. Specifically, visual evaluation confirms both models struggle significantly more at the mid-wings 
while maintaining higher accuracy near the center and the extreme baseline.
\subsection{Adjusting for true frequencies}
The original dataset has roughly equal representation for each competition, but in reality, the NBA accounts for 60\% of shots while the other four types account for only 10\% each. To estimate real-world performance, we calculated the accuracy and log-loss for each competition separately and then used a weighted average based on these true percentages. This allows us to see how the models perform on the actual data distribution without retraining.
\begin{table}[htpb]
    \centering
    \caption{Model Performance Adjusted for True Competition Frequencies}
    \label{tab:reweighted_metrics}
    \begin{tabular}{lcccc}
        \toprule
        & \multicolumn{2}{c}{\textbf{Unweighted (Original)}} & \multicolumn{2}{c}{\textbf{Weighted (True Freq.)}} \\
        \cmidrule(lr){2-3} \cmidrule(lr){4-5}
        \textbf{Model} & \textbf{Accuracy} & \textbf{Log-Loss} & \textbf{Accuracy} & \textbf{Log-Loss} \\
        \midrule
        Decision Tree (Nested CV) & 0.749 & 0.844 & 0.753 & 0.840 \\
        Logistic Regression & 0.736 & 0.673 & \textbf{0.753} & \textbf{0.656} \\
        \bottomrule
    \end{tabular}
\end{table}
As shown in Table \ref{tab:reweighted_metrics}, both models perform better when adjusted for the true frequencies. The accuracy for both increases to 0.753, and the log-loss for Logistic Regression improves. This suggests the models are more accurate at predicting NBA shots than shots from other competitions. Since the NBA makes up a much larger portion of the real-world data, the models' overall reliability is actually higher than the initial results suggested.
\end{document}